% ============================================================================
% Probe-IPA Dokumentation – Lion Hereqi
% Full-Stack Bibliotheksverwaltung
% Datum: 09.03.2026
% ============================================================================
\documentclass[a4paper,12pt,twoside]{report}

% --- Pakete ---
\usepackage[utf8]{inputenc}
\usepackage[T1]{fontenc}
\usepackage[ngerman]{babel}
\usepackage{geometry}
\geometry{a4paper, left=2.5cm, right=2.5cm, top=2.5cm, bottom=2.5cm}
\usepackage{graphicx}
\usepackage{xcolor}
\usepackage{hyperref}
\usepackage{listings}
\usepackage{booktabs}
\usepackage{tabularx}
\usepackage{longtable}
\usepackage{array}
\usepackage{fancyhdr}
\usepackage{amssymb}
\usepackage{caption}
\usepackage{float}
\usepackage{parskip}
\usepackage{textcomp}
\usepackage{pdflscape}
\usepackage{tikz}

% --- Farben (ti&m Design System) ---
\definecolor{timprimary}{HTML}{00918F}
\definecolor{timsecondary}{HTML}{002650}
\definecolor{timtext}{HTML}{202020}
\definecolor{timerror}{HTML}{DC2918}
\definecolor{codebg}{HTML}{F5F5F5}
\definecolor{codeframe}{HTML}{E0E0E0}

% --- Hyperref Einstellungen ---
\hypersetup{
    colorlinks=true,
    linkcolor=timsecondary,
    urlcolor=timprimary,
    citecolor=timprimary,
    pdfauthor={Lion Hereqi},
    pdftitle={Probe-IPA 2026 – Full-Stack Bibliotheksverwaltung},
}

% --- Code-Listings ---
\lstdefinelanguage{java}{
    keywords={package, import, public, private, protected, class, interface, extends, implements, static, final, void, return, if, else, for, while, new, this, super, true, false, null, try, catch, throw, throws},
    keywordstyle=\color{timsecondary}\bfseries,
    ndkeywords={String, Long, List, Book, Loan, Entity, Table, Id, Column, Service, Repository, Controller, GetMapping, PostMapping, PutMapping, DeleteMapping, RequestBody, PathVariable, ResponseEntity, Autowired},
    ndkeywordstyle=\color{timprimary}\bfseries,
    sensitive=true,
    comment=[l]{//},
    morecomment=[s]{/*}{*/},
    commentstyle=\color{gray}\itshape,
    string=[b]",
    stringstyle=\color{timerror},
}

\lstdefinelanguage{tsx}{
    keywords={function, const, let, var, return, if, else, for, while, import, export, from, async, await, try, catch, throw, new, typeof, instanceof, default, interface, type, extends, implements, class, this, super, true, false, null, undefined},
    keywordstyle=\color{timsecondary}\bfseries,
    ndkeywords={useState, useEffect, useKeycloak, React, Book, Loan},
    ndkeywordstyle=\color{timprimary}\bfseries,
    sensitive=true,
    comment=[l]{//},
    morecomment=[s]{/*}{*/},
    commentstyle=\color{gray}\itshape,
    string=[b]",
    morestring=[b]',
    morestring=[b]`,
    stringstyle=\color{timerror},
}

\lstset{
    basicstyle=\ttfamily\footnotesize,
    backgroundcolor=\color{codebg},
    frame=single,
    rulecolor=\color{codeframe},
    breaklines=true,
    breakatwhitespace=false,
    showstringspaces=false,
    tabsize=2,
    numbers=left,
    numberstyle=\tiny\color{gray},
    numbersep=8pt,
    xleftmargin=15pt,
    framexleftmargin=15pt,
    aboveskip=10pt,
    belowskip=10pt,
    captionpos=b,
    literate={ä}{{\"a}}1 {ö}{{\"o}}1 {ü}{{\"u}}1 {Ä}{{\"A}}1 {Ö}{{\"O}}1 {Ü}{{\"U}}1 {ß}{{\ss}}1,
}

% --- Kopf-/Fusszeile ---
\pagestyle{fancy}
\fancyhf{}
\fancyhead[L]{\small Lion Hereqi -- Probe-IPA 2026}
\fancyhead[R]{\small Full-Stack Bibliotheksverwaltung}
\fancyfoot[C]{\thepage}
\renewcommand{\headrulewidth}{0.4pt}
\renewcommand{\footrulewidth}{0pt}

% --- Checkmark / Crossmark ---
\newcommand{\cmark}{\textcolor{green!60!black}{$\checkmark$}}
\newcommand{\xmark}{\textcolor{red}{$\times$}}

% --- Tabellentypen ---
\newcolumntype{L}[1]{>{\raggedright\arraybackslash}p{#1}}
\newcolumntype{R}[1]{>{\raggedleft\arraybackslash}p{#1}}
\newcolumntype{C}[1]{>{\centering\arraybackslash}p{#1}}

% ============================================================================
\begin{document}

% --- Titelseite ---
\begin{titlepage}
    \centering
    \vspace*{1cm}

    {\Large\bfseries ti\&m AG}\\[0.3cm]
    {\small Zürich \textbar\ Bern \textbar\ Basel \textbar\ Singapur \textbar\ Frankfurt}\\[0.2cm]
    {\small Buckhauserstrasse 24, 8048 Zürich}\\[0.2cm]
    {\small T +41 44 497 75 00 \quad F +41 44 497 75 01}\\[0.2cm]
    {\small \href{mailto:info@ti8m.com}{info@ti8m.com} \quad \href{https://www.ti8m.com}{www.ti8m.com}}\\[2cm]

    \rule{\textwidth}{1pt}\\[1.5cm]

    {\Huge\bfseries\color{timsecondary} Lion Hereqi}\\[0.5cm]
    {\LARGE\bfseries Probe-IPA 05.03. -- 09.03.2026}\\[1.5cm]

    {\huge\bfseries\color{timprimary} Full-Stack\\Bibliotheksverwaltung}\\[2cm]

    \rule{\textwidth}{1pt}\\[1cm]

    {\large Autor: Lion Hereqi}\\[0.3cm]
    {\large Zuletzt geändert am: 09.03.2026}\\

    \vfill
\end{titlepage}

% --- Inhaltsverzeichnis ---
\tableofcontents
\newpage

% ============================================================================
% TEIL 1 – UMFELD UND ABLAUF
% ============================================================================
\chapter{Teil 1 -- Umfeld und Ablauf}

% ----------------------------------------------------------------------------
\section{Vorwort}
Das Ziel dieser Probe-IPA ist es, meine Fachkenntnisse als Applikationsentwickler unter Beweis zu stellen. Im Rahmen dieser Arbeit wurde eine webbasierte Full-Stack-Applikation zur Verwaltung von Büchern in einer Bibliothek entwickelt.

Das Dokument gliedert sich in zwei Teile:
\begin{itemize}
    \item \textbf{Teil 1:} Umfeld und Ablauf -- Organisatorische Aspekte der Arbeit
    \item \textbf{Teil 2:} Projekt und Umsetzung -- Technische Dokumentation
\end{itemize}

\newpage

% ----------------------------------------------------------------------------
\section{Aufgabenstellung}

\subsection{Titel der Arbeit}
Full-Stack Bibliotheksverwaltung

\subsection{Ausgangslage}
In einer Bibliothek werden Bücher aktuell nicht digital, sondern noch mit einem alten Akten-Order-System auf Papier verwaltet. Dadurch entsteht ein deutlicher Mehraufwand für die Angestellten der Bibliothek.

Ziel dieser Probearbeit ist die Entwicklung eines webbasierten Full-Stack-Tools zur strukturierten Erfassung und Verwaltung von Büchern.

\subsection{Detaillierte Aufgabenstellung}
Das System wird als moderne Webanwendung umgesetzt mit folgenden Technologien:
\begin{itemize}
    \item Frontend: React + Vite + TypeScript + Shadcn UI
    \item Backend: Java + Spring Boot
    \item Datenbank: MariaDB
    \item Authentifizierung \& Rollenverwaltung: Keycloak via OIDC
\end{itemize}

Die Kernfunktionen, die diese Probe-IPA lösen soll:
\begin{itemize}
    \item Erfassen von Büchern in der Bibliothek
    \item Zuordnen von Büchern als Verliehen/Verfügbar
    \item Wenn ein Buch verliehen ist, muss es einem Nutzer zugeordnet sein
    \item Nutzer können ihre ausgeliehenen Bücher auflisten
\end{itemize}

\subsection{Rollen- und Sicherheitskonzept}
Das System unterscheidet zwischen zwei Rollen:
\begin{enumerate}
    \item \textbf{Mitarbeiter (STAFF):} Kann neue Bücher anlegen, Bücher an Kunden verleihen und als zurückgegeben markieren.
    \item \textbf{Kunde (CUSTOMER):} Kann sich die Bücher ansehen, die ihm aktuell zugewiesen sind.
\end{enumerate}

\subsection{Abgrenzung}
\begin{itemize}
    \item Die Applikation ist eine Webanwendung, mobile Unterstützung ist nicht erforderlich.
    \item Benutzerverwaltung erfolgt über Keycloak, nicht in der Applikation selbst.
    \item Es werden keine externen Schnittstellen zu anderen Systemen implementiert.
\end{itemize}

\newpage

% ----------------------------------------------------------------------------
\section{Mittel und Methoden}

\subsection{Hardware}

\begin{table}[H]
\centering
\caption{Verwendete Hardware}
\label{tab:hardware}
\begin{tabularx}{\textwidth}{l X}
\toprule
\textbf{Eigenschaft} & \textbf{Beschreibung} \\
\midrule
Marke & Apple \\
Modell & MacBook Pro 16", Nov. 2024 \\
Chip & Apple M4 Max \\
Arbeitsspeicher & 64 GB \\
Betriebssystem & macOS Sequoia 15.7.2 \\
\bottomrule
\end{tabularx}
\end{table}

\subsection{Arbeitsplatz}
Die gesamte Durchführung der Probe-IPA erfolgt vor Ort bei der ti\&m AG in Zürich.

\subsection{Software}

\begin{table}[H]
\centering
\caption{Verwendete Software}
\label{tab:software}
\begin{tabularx}{\textwidth}{l X}
\toprule
\textbf{Tool} & \textbf{Beschreibung} \\
\midrule
Cursor IDE & Entwicklungsumgebung mit KI-Unterstützung \\
Docker Desktop & Container-Verwaltung für MariaDB und Keycloak \\
Git & Versionskontrolle \\
GitHub & Remote Repository (\href{https://github.com/hereqi/Probe-IPA-Buecherrei}{github.com/hereqi/Probe-IPA-Buecherrei}) \\
Chrome & Browser zum Testen \\
Postman & API-Testing \\
\bottomrule
\end{tabularx}
\end{table}

\subsection{Technologien}

\begin{table}[H]
\centering
\caption{Verwendete Technologien}
\begin{tabularx}{\textwidth}{l l X}
\toprule
\textbf{Bereich} & \textbf{Technologie} & \textbf{Version/Beschreibung} \\
\midrule
Frontend & React & 18.x mit TypeScript \\
Frontend & Vite & Build-Tool \\
Frontend & Tailwind CSS & Utility-First CSS Framework \\
Frontend & Shadcn UI & Komponentenbibliothek \\
Backend & Java & 17 \\
Backend & Spring Boot & 3.5.x \\
Backend & Spring Security & OAuth2 Resource Server \\
Backend & Spring Data JPA & Hibernate ORM \\
Datenbank & MariaDB & 11.x \\
Auth & Keycloak & 23.x (OIDC) \\
Infra & Docker Compose & Container-Orchestrierung \\
\bottomrule
\end{tabularx}
\end{table}

\newpage

% ----------------------------------------------------------------------------
\section{Neue Lerninhalte}

Im Rahmen dieser Probe-IPA wurden folgende neue Kenntnisse erworben:
\begin{itemize}
    \item Integration von Keycloak mit Spring Boot (OAuth2 Resource Server)
    \item Rollenbasierte Zugriffskontrolle mit JWT-Token
    \item Shadcn UI Komponentenbibliothek
    \item Docker Compose für Multi-Container-Umgebungen
    \item Full-Stack Architektur mit React und Spring Boot
\end{itemize}

% ----------------------------------------------------------------------------
\section{Vorkenntnisse}

\subsection{Technologien}
\begin{itemize}
    \item Entwicklung mit React und TypeScript
    \item Entwicklung mit Java und Spring Boot
    \item Styling mit Tailwind CSS
    \item Versionierung mit Git
\end{itemize}

\subsection{Software und Tools}
\begin{itemize}
    \item Cursor IDE / VS Code
    \item Git und GitHub
    \item Docker
\end{itemize}

\newpage

% ----------------------------------------------------------------------------
\section{Projektaufbauorganisation}

\subsection{Organigramm}

\begin{figure}[H]
\centering
\begin{tikzpicture}[
    node distance=1.5cm and 2cm,
    every node/.style={draw, rounded corners, minimum width=4cm, minimum height=1cm, align=center, font=\small}
]
    \node (kandidat) {Kandidat\\Lion Hereqi};
    \node (vf) [above=1.5cm of kandidat] {Fachvorgesetzter\\Georg Steinmetz};
    \node (he) [above right=1cm and 2.5cm of kandidat] {Hauptexperte/in\\[PLATZHALTER]};
    \node (ne) [below right=1cm and 2.5cm of kandidat] {Nebenexperte/in\\[PLATZHALTER]};
    \draw[->] (vf) -- (kandidat);
    \draw[->] (he) -- (kandidat);
    \draw[->] (ne) -- (kandidat);
\end{tikzpicture}
\caption{Organigramm}
\label{fig:organigramm}
\end{figure}

\subsection{Kontaktdaten}

\textbf{Betrieb (Durchführungsort)}\\
ti\&m AG\\
Buckhauserstrasse 24\\
8048 Zürich

\textbf{Kandidat}\\
Lion Hereqi\\
[PLATZHALTER: Kontaktdaten]

\textbf{Fachvorgesetzter}\\
Georg Steinmetz\\
ti\&m AG

\newpage

% ----------------------------------------------------------------------------
\section{Versionierung}

\subsection{Code-Versionierung}
Der gesamte Quellcode wird mit Git versioniert und auf GitHub gespeichert:

\begin{center}
\url{https://github.com/hereqi/Probe-IPA-Buecherrei}
\end{center}

\subsection{Dokumentations-Versionierung}
Die Dokumentation wird ebenfalls im Git-Repository unter \texttt{/docs} versioniert.

\newpage

% ----------------------------------------------------------------------------
\section{Zeitplan}

% Querformat-Seite für Zeitplan
\begin{landscape}
\thispagestyle{empty}

\vspace*{1cm}
{\Large\bfseries Zeitplan -- Probe-IPA Bibliotheksverwaltung}\\[0.5cm]
{\normalsize Zeitraum: 05.03.2026 -- 09.03.2026 (5 Tage)}\\[1cm]

\begin{center}
\textit{[PLATZHALTER: Zeitplan als Gantt-Diagramm oder Tabelle einfügen]}\\[2cm]

\fbox{\parbox{0.9\linewidth}{\centering\vspace{8cm}
\textbf{PLATZHALTER}\\[0.5cm]
Hier wird der detaillierte Zeitplan eingefügt.\\
Format: Gantt-Diagramm oder tabellarische Übersicht\\
mit Tätigkeiten, Soll- und Ist-Zeiten.
\vspace{8cm}}}
\end{center}

\end{landscape}

\newpage

% ============================================================================
% TEIL 2 – PROJEKT UND UMSETZUNG
% ============================================================================
\chapter{Teil 2 -- Projekt und Umsetzung}

% ----------------------------------------------------------------------------
\section{Anforderungsanalyse}

\subsection{Funktionale Anforderungen}

\begin{table}[H]
\centering
\caption{Funktionale Anforderungen}
\label{tab:fa}
\begin{tabularx}{\textwidth}{c c X c}
\toprule
\textbf{ID} & \textbf{Prio} & \textbf{Beschreibung} & \textbf{Status} \\
\midrule
FA-01 & Muss & Bücher erfassen (Titel, Autor, ISBN) & \cmark \\
FA-02 & Muss & Bücher an Kunden ausleihen & \cmark \\
FA-03 & Muss & Ausgeliehene Bücher zurückgeben & \cmark \\
FA-04 & Muss & Kunde sieht eigene aktive Ausleihen & \cmark \\
FA-05 & Muss & Rollenbasierte Zugriffskontrolle (STAFF/CUSTOMER) & \cmark \\
FA-06 & Soll & Übersicht aller Bücher mit Status & \cmark \\
FA-07 & Soll & Übersicht aller aktiven Ausleihen (Mitarbeiter) & \cmark \\
\bottomrule
\end{tabularx}
\end{table}

\subsection{Nicht-funktionale Anforderungen}

\begin{table}[H]
\centering
\caption{Nicht-funktionale Anforderungen}
\label{tab:nfa}
\begin{tabularx}{\textwidth}{c l X c}
\toprule
\textbf{ID} & \textbf{Kategorie} & \textbf{Beschreibung} & \textbf{Status} \\
\midrule
NFA-01 & Sicherheit & Authentifizierung via Keycloak (OIDC) & \cmark \\
NFA-02 & Sicherheit & JWT-Token Validierung im Backend & \cmark \\
NFA-03 & Sicherheit & Rollenbasierte API-Absicherung & \cmark \\
NFA-04 & Usability & Konsistente Benutzeroberfläche & \cmark \\
NFA-05 & Wartbarkeit & Strukturierter, dokumentierter Code & \cmark \\
NFA-06 & Betrieb & Reproduzierbare Entwicklungsumgebung (Docker) & \cmark \\
\bottomrule
\end{tabularx}
\end{table}

\newpage

% ----------------------------------------------------------------------------
\section{Use Cases}

\subsection{Use Case Diagramm}

\begin{figure}[H]
\centering
\fbox{\parbox{0.85\textwidth}{\centering\vspace{4cm}
\textbf{PLATZHALTER}\\[0.5cm]
Use Case Diagramm hier einfügen.\\
Datei: \texttt{screenshots/use-case-diagram.png}
\vspace{4cm}}}
\caption{Use Case Diagramm}
\label{fig:usecase}
\end{figure}

\newpage

\subsection{Use Case Beschreibungen}

\subsubsection{UC-01: Anmelden}
\begin{table}[H]
\begin{tabularx}{\textwidth}{l X}
\toprule
\textbf{Use Case} & UC-01: Anmelden \\
\midrule
Akteure & Mitarbeiter, Kunde \\
Vorbedingung & Benutzer hat einen Keycloak-Account \\
Ablauf & 1. Benutzer klickt auf ``Anmelden''\\
& 2. Weiterleitung zur Keycloak Login-Seite\\
& 3. Eingabe von Benutzername und Passwort\\
& 4. Weiterleitung zurück zur Applikation \\
Nachbedingung & Benutzer ist eingeloggt, JWT-Token erhalten \\
\bottomrule
\end{tabularx}
\end{table}

\subsubsection{UC-02: Buch erfassen}
\begin{table}[H]
\begin{tabularx}{\textwidth}{l X}
\toprule
\textbf{Use Case} & UC-02: Buch erfassen \\
\midrule
Akteure & Mitarbeiter \\
Vorbedingung & Mitarbeiter ist angemeldet (Rolle: STAFF) \\
Ablauf & 1. Mitarbeiter öffnet den Mitarbeiter-Bereich\\
& 2. Gibt Titel, Autor und optional ISBN ein\\
& 3. Klickt auf ``Hinzufügen'' \\
Nachbedingung & Buch ist in der Datenbank gespeichert mit Status ``AVAILABLE'' \\
\bottomrule
\end{tabularx}
\end{table}

\subsubsection{UC-03: Buch ausleihen}
\begin{table}[H]
\begin{tabularx}{\textwidth}{l X}
\toprule
\textbf{Use Case} & UC-03: Buch ausleihen \\
\midrule
Akteure & Mitarbeiter \\
Vorbedingung & Buch ist verfügbar (Status: AVAILABLE) \\
Ablauf & 1. Mitarbeiter wählt ein verfügbares Buch aus dem Dropdown\\
& 2. Gibt den Kundennamen ein\\
& 3. Klickt auf ``Ausleihen'' \\
Nachbedingung & Buch-Status ist ``BORROWED'', Ausleihe ist erstellt \\
\bottomrule
\end{tabularx}
\end{table}

\subsubsection{UC-04: Buch zurückgeben}
\begin{table}[H]
\begin{tabularx}{\textwidth}{l X}
\toprule
\textbf{Use Case} & UC-04: Buch zurückgeben \\
\midrule
Akteure & Mitarbeiter \\
Vorbedingung & Ausleihe existiert mit Status ``ACTIVE'' \\
Ablauf & 1. Mitarbeiter sieht Liste der aktiven Ausleihen\\
& 2. Klickt bei einer Ausleihe auf ``Zurückgeben'' \\
Nachbedingung & Ausleihe-Status ist ``RETURNED'', Buch-Status ist ``AVAILABLE'' \\
\bottomrule
\end{tabularx}
\end{table}

\subsubsection{UC-05: Eigene Ausleihen anzeigen}
\begin{table}[H]
\begin{tabularx}{\textwidth}{l X}
\toprule
\textbf{Use Case} & UC-05: Eigene Ausleihen anzeigen \\
\midrule
Akteure & Kunde \\
Vorbedingung & Kunde ist angemeldet (Rolle: CUSTOMER) \\
Ablauf & 1. Kunde öffnet ``Meine Ausleihen''\\
& 2. Liste der aktiven Ausleihen wird angezeigt \\
Nachbedingung & -- \\
\bottomrule
\end{tabularx}
\end{table}

\newpage

% ----------------------------------------------------------------------------
\section{Datenmodell}

\subsection{ER-Diagramm}

\begin{figure}[H]
\centering
\fbox{\parbox{0.85\textwidth}{\centering\vspace{4cm}
\textbf{PLATZHALTER}\\[0.5cm]
ER-Diagramm hier einfügen.\\
Datei: \texttt{screenshots/er-diagram.png}
\vspace{4cm}}}
\caption{ER-Diagramm}
\label{fig:er}
\end{figure}

\newpage

\subsection{Entität: Book}

\begin{table}[H]
\centering
\caption{Entität: Book (Buch)}
\begin{tabularx}{\textwidth}{l l l X}
\toprule
\textbf{Attribut} & \textbf{Typ} & \textbf{Constraint} & \textbf{Beschreibung} \\
\midrule
id & BIGINT & PK, Auto & Primärschlüssel \\
title & VARCHAR(255) & NOT NULL & Buchtitel \\
author & VARCHAR(255) & NOT NULL & Autor des Buches \\
isbn & VARCHAR(50) & UNIQUE & ISBN-Nummer \\
status & VARCHAR(50) & -- & AVAILABLE oder BORROWED \\
\bottomrule
\end{tabularx}
\end{table}

\subsection{Entität: Loan}

\begin{table}[H]
\centering
\caption{Entität: Loan (Ausleihe)}
\begin{tabularx}{\textwidth}{l l l X}
\toprule
\textbf{Attribut} & \textbf{Typ} & \textbf{Constraint} & \textbf{Beschreibung} \\
\midrule
id & BIGINT & PK, Auto & Primärschlüssel \\
book\_id & BIGINT & FK, NOT NULL & Fremdschlüssel zu Book \\
userId & VARCHAR(255) & NOT NULL & Keycloak User-ID \\
userName & VARCHAR(255) & NOT NULL & Benutzername \\
loanDate & DATE & NOT NULL & Ausleihdatum \\
dueDate & DATE & NOT NULL & Fälligkeitsdatum (4 Wochen) \\
returnDate & DATE & -- & Rückgabedatum (nullable) \\
status & VARCHAR(50) & -- & ACTIVE oder RETURNED \\
\bottomrule
\end{tabularx}
\end{table}

\newpage

% ----------------------------------------------------------------------------
\section{Architektur}

\subsection{Systemübersicht}

\begin{figure}[H]
\centering
\fbox{\parbox{0.85\textwidth}{\centering\vspace{4cm}
\textbf{PLATZHALTER}\\[0.5cm]
Architektur-Diagramm hier einfügen.\\
Datei: \texttt{screenshots/architecture.png}
\vspace{4cm}}}
\caption{Systemarchitektur}
\label{fig:architektur}
\end{figure}

\subsection{Schichtenmodell}

\begin{verbatim}
+-------------------------------------------------------------+
|                        Frontend                             |
|               React + TypeScript + Shadcn UI                |
|                    (Port 5173)                              |
+-------------------------------------------------------------+
                            |
                            | HTTP/REST + JWT
                            v
+-------------------------------------------------------------+
|                        Backend                              |
|                   Spring Boot REST API                      |
|                     (Port 8080)                             |
|  +-------------+  +-------------+  +-------------------+    |
|  | Controller  |  |   Service   |  |    Repository     |    |
|  +-------------+  +-------------+  +-------------------+    |
+-------------------------------------------------------------+
                            |
                            | JPA/Hibernate
                            v
+-------------------------------------------------------------+
|                       Datenbank                             |
|                        MariaDB                              |
|                      (Port 3306)                            |
+-------------------------------------------------------------+

+-------------------------------------------------------------+
|                   Authentifizierung                         |
|                       Keycloak                              |
|                      (Port 8180)                            |
+-------------------------------------------------------------+
\end{verbatim}

\newpage

\subsection{Projektstruktur}

\textbf{Backend (/backend):}
\begin{verbatim}
src/main/java/ch/library/backend/
├── BackendApplication.java      # Spring Boot Einstiegspunkt
├── config/
│   └── SecurityConfig.java      # OAuth2/JWT Konfiguration
├── controller/
│   ├── BookController.java      # REST Endpoints für Bücher
│   └── LoanController.java      # REST Endpoints für Ausleihen
├── dto/
│   └── BookRequest.java         # Request DTO mit Validierung
├── entity/
│   ├── Book.java                # JPA Entity Buch
│   └── Loan.java                # JPA Entity Ausleihe
├── exception/
│   ├── GlobalExceptionHandler.java
│   ├── ResourceNotFoundException.java
│   └── BadRequestException.java
├── repository/
│   ├── BookRepository.java
│   └── LoanRepository.java
└── service/
    ├── BookService.java
    └── LoanService.java
\end{verbatim}

\textbf{Frontend (/frontend):}
\begin{verbatim}
src/
├── App.tsx                      # Hauptkomponente mit Routing
├── main.tsx                     # Einstiegspunkt
├── index.css                    # Tailwind CSS
├── components/
│   └── ui/                      # Shadcn UI Komponenten
├── pages/
│   ├── BooksPage.tsx            # Bücherübersicht
│   ├── MyLoansPage.tsx          # Meine Ausleihen (Kunde)
│   └── StaffPage.tsx            # Mitarbeiter-Bereich
├── services/
│   ├── api.ts                   # Axios API Client
│   └── keycloak.ts              # Keycloak Konfiguration
└── types/
    └── index.ts                 # TypeScript Interfaces
\end{verbatim}

\newpage

% ----------------------------------------------------------------------------
\section{REST API}

\subsection{Endpoints}

\begin{table}[H]
\centering
\caption{REST API Endpoints}
\begin{tabularx}{\textwidth}{l l X l}
\toprule
\textbf{Methode} & \textbf{Pfad} & \textbf{Beschreibung} & \textbf{Rolle} \\
\midrule
GET & /api/books & Alle Bücher abrufen & Alle \\
GET & /api/books/\{id\} & Einzelnes Buch abrufen & Alle \\
GET & /api/books/available & Verfügbare Bücher & Alle \\
POST & /api/books & Neues Buch erstellen & STAFF \\
PUT & /api/books/\{id\} & Buch aktualisieren & STAFF \\
DELETE & /api/books/\{id\} & Buch löschen & STAFF \\
\midrule
GET & /api/loans & Alle Ausleihen & STAFF \\
GET & /api/loans/my & Eigene Ausleihen & Auth \\
POST & /api/loans & Neue Ausleihe erstellen & STAFF \\
PUT & /api/loans/\{id\}/return & Buch zurückgeben & STAFF \\
\bottomrule
\end{tabularx}
\end{table}

\newpage

% ----------------------------------------------------------------------------
\section{Implementierung}

\subsection{Backend: Entity Book}

\begin{lstlisting}[language=java, caption=Book.java]
@Entity
@Table(name = "books")
public class Book {

    @Id
    @GeneratedValue(strategy = GenerationType.IDENTITY)
    private Long id;

    @Column(nullable = false)
    private String title;

    @Column(nullable = false)
    private String author;

    @Column(unique = true)
    private String isbn;

    private String status = "AVAILABLE";

    // Getter und Setter
}
\end{lstlisting}

\subsection{Backend: Security Config}

\begin{lstlisting}[language=java, caption=SecurityConfig.java (Auszug)]
@Configuration
@EnableWebSecurity
public class SecurityConfig {

    @Bean
    public SecurityFilterChain filterChain(HttpSecurity http) throws Exception {
        http
            .cors(cors -> {})
            .csrf(csrf -> csrf.disable())
            .authorizeHttpRequests(auth -> auth
                .requestMatchers(HttpMethod.GET, "/api/books/**").permitAll()
                .requestMatchers(HttpMethod.POST, "/api/books/**").hasRole("STAFF")
                .requestMatchers("/api/loans/my").authenticated()
                .requestMatchers("/api/loans/**").hasRole("STAFF")
                .anyRequest().authenticated()
            )
            .oauth2ResourceServer(oauth2 -> oauth2
                .jwt(jwt -> jwt.jwtAuthenticationConverter(jwtAuthenticationConverter()))
            );
        return http.build();
    }
}
\end{lstlisting}

\newpage

\subsection{Frontend: BooksPage}

\begin{lstlisting}[language=tsx, caption=BooksPage.tsx (Auszug)]
export default function BooksPage() {
  const [books, setBooks] = useState<Book[]>([]);
  const [loading, setLoading] = useState(true);

  useEffect(() => {
    bookApi.getAll()
      .then(setBooks)
      .finally(() => setLoading(false));
  }, []);

  if (loading) return <p>Laden...</p>;

  return (
    <Card>
      <CardHeader>
        <CardTitle>Buecher</CardTitle>
      </CardHeader>
      <CardContent>
        <Table>
          {/* Tabelle mit Buechern */}
        </Table>
      </CardContent>
    </Card>
  );
}
\end{lstlisting}

\newpage

% ----------------------------------------------------------------------------
\section{Sicherheit}

\subsection{Authentifizierung}
Die Authentifizierung erfolgt über Keycloak mit dem OpenID Connect (OIDC) Protokoll:

\begin{enumerate}
    \item Benutzer klickt auf ``Anmelden''
    \item Weiterleitung zur Keycloak Login-Seite
    \item Nach erfolgreichem Login erhält das Frontend ein JWT-Token
    \item Token wird bei jedem API-Request im Authorization-Header mitgeschickt
    \item Backend validiert das Token gegen Keycloak
\end{enumerate}

\subsection{Autorisierung}
Rollenbasierte Zugriffskontrolle (RBAC):
\begin{itemize}
    \item \textbf{STAFF:} Vollzugriff auf alle Funktionen
    \item \textbf{CUSTOMER:} Nur lesender Zugriff auf eigene Daten
\end{itemize}

\subsection{Implementierte Sicherheitsmassnahmen}
\begin{itemize}
    \item JWT-Token Validierung im Backend
    \item Keycloak Realm-Rollen werden korrekt gemappt
    \item CORS-Konfiguration für Frontend-Origin
    \item CSRF-Schutz (stateless durch JWT)
    \item SQL-Injection-Schutz durch JPA/Hibernate (Prepared Statements)
    \item Eingabevalidierung mit Jakarta Validation (@NotBlank, @Size)
\end{itemize}

\newpage

% ----------------------------------------------------------------------------
\section{Fehlerbehandlung und Logging}

\subsection{Fehlerbehandlung}
Die Applikation verwendet einen globalen Exception Handler:

\begin{lstlisting}[language=java, caption=GlobalExceptionHandler.java (Auszug)]
@RestControllerAdvice
public class GlobalExceptionHandler {

    private static final Logger logger = 
        LoggerFactory.getLogger(GlobalExceptionHandler.class);

    @ExceptionHandler(ResourceNotFoundException.class)
    public ResponseEntity<Map<String, Object>> handleNotFound(
            ResourceNotFoundException ex) {
        logger.warn("Resource nicht gefunden: {}", ex.getMessage());
        return buildResponse(HttpStatus.NOT_FOUND, ex.getMessage());
    }

    @ExceptionHandler(MethodArgumentNotValidException.class)
    public ResponseEntity<Map<String, Object>> handleValidation(
            MethodArgumentNotValidException ex) {
        // Validierungsfehler verarbeiten
    }
}
\end{lstlisting}

\subsection{Logging}
Logging erfolgt mit SLF4J/Logback:
\begin{itemize}
    \item INFO: Normale Operationen (Buch erstellt, Ausleihe erstellt)
    \item WARN: Erwartete Fehler (Resource nicht gefunden)
    \item ERROR: Unerwartete Fehler (Exceptions)
\end{itemize}

\newpage

% ----------------------------------------------------------------------------
\section{Testing}

\subsection{Testfälle}

\begin{table}[H]
\centering
\caption{Durchgeführte Testfälle}
\begin{tabularx}{\textwidth}{c X X c}
\toprule
\textbf{ID} & \textbf{Beschreibung} & \textbf{Erwartetes Ergebnis} & \textbf{OK} \\
\midrule
T-01 & Login als Mitarbeiter & Zugriff auf Staff-Bereich & \cmark \\
T-02 & Login als Kunde & Nur ``Meine Ausleihen'' sichtbar & \cmark \\
T-03 & Buch erstellen (Mitarbeiter) & Buch erscheint in Liste & \cmark \\
T-04 & Buch erstellen ohne Titel & Validierungsfehler & \cmark \\
T-05 & Buch ausleihen & Status wechselt zu ``Ausgeliehen'' & \cmark \\
T-06 & Buch zurückgeben & Status wechselt zu ``Verfügbar'' & \cmark \\
T-07 & Kunde sieht eigene Ausleihen & Liste wird angezeigt & \cmark \\
T-08 & Kunde versucht Buch zu erstellen & Kein Zugriff (403) & \cmark \\
T-09 & Nicht angemeldet: Staff-Bereich & Redirect zu Login & \cmark \\
\bottomrule
\end{tabularx}
\end{table}

\newpage

% ----------------------------------------------------------------------------
\section{Screenshots}

\subsection{Bücher-Übersicht}
\begin{figure}[H]
\centering
\fbox{\parbox{0.85\textwidth}{\centering\vspace{5cm}
\textbf{PLATZHALTER}\\[0.5cm]
Screenshot der Bücher-Seite hier einfügen.\\
Datei: \texttt{screenshots/books-page.png}
\vspace{5cm}}}
\caption{Bücher-Übersicht}
\end{figure}

\newpage

\subsection{Mitarbeiter-Bereich}
\begin{figure}[H]
\centering
\fbox{\parbox{0.85\textwidth}{\centering\vspace{5cm}
\textbf{PLATZHALTER}\\[0.5cm]
Screenshot des Mitarbeiter-Bereichs hier einfügen.\\
Datei: \texttt{screenshots/staff-page.png}
\vspace{5cm}}}
\caption{Mitarbeiter-Bereich}
\end{figure}

\newpage

\subsection{Meine Ausleihen (Kundenansicht)}
\begin{figure}[H]
\centering
\fbox{\parbox{0.85\textwidth}{\centering\vspace{5cm}
\textbf{PLATZHALTER}\\[0.5cm]
Screenshot der Kundenansicht hier einfügen.\\
Datei: \texttt{screenshots/my-loans-page.png}
\vspace{5cm}}}
\caption{Meine Ausleihen}
\end{figure}

\newpage

\subsection{Keycloak Login}
\begin{figure}[H]
\centering
\fbox{\parbox{0.85\textwidth}{\centering\vspace{5cm}
\textbf{PLATZHALTER}\\[0.5cm]
Screenshot der Keycloak Login-Seite hier einfügen.\\
Datei: \texttt{screenshots/keycloak-login.png}
\vspace{5cm}}}
\caption{Keycloak Login}
\end{figure}

\newpage

% ----------------------------------------------------------------------------
\section{Reflexion}

\subsection{Zielerreichung}
Alle funktionalen Anforderungen wurden erfolgreich umgesetzt:
\begin{itemize}
    \item[\cmark] Bücher können erfasst werden
    \item[\cmark] Ausleihen können erstellt und zurückgegeben werden
    \item[\cmark] Rollenbasierte Navigation funktioniert korrekt
    \item[\cmark] Kunden sehen nur ihre eigenen Ausleihen
    \item[\cmark] Authentifizierung über Keycloak funktioniert
\end{itemize}

\subsection{Aufgetretene Probleme und Lösungen}

\begin{table}[H]
\centering
\caption{Aufgetretene Probleme}
\begin{tabularx}{\textwidth}{X X}
\toprule
\textbf{Problem} & \textbf{Lösung} \\
\midrule
Keycloak Rollen wurden im Backend nicht erkannt & KeycloakRoleConverter implementiert, der realm\_access.roles ausliest \\
\midrule
Tailwind CSS Konflikt mit Shadcn UI & CSS-Variablen korrekt konfiguriert, @layer Direktiven entfernt \\
\midrule
``Meine Ausleihen'' zeigte keine Daten & Suche nach userName statt userId, da Keycloak-ID nicht direkt verfügbar \\
\midrule
Vite Import-Alias funktionierte nicht & tsconfig.json und vite.config.ts korrekt konfiguriert \\
\bottomrule
\end{tabularx}
\end{table}

\subsection{Neue Erkenntnisse}
\begin{itemize}
    \item Full-Stack Entwicklung mit React und Spring Boot
    \item Integration von Keycloak für Authentifizierung und Autorisierung
    \item Docker Compose für lokale Entwicklungsumgebung
    \item Shadcn UI als moderne Komponentenbibliothek
\end{itemize}

\subsection{Fazit}
Die Probe-IPA wurde erfolgreich abgeschlossen. Das entwickelte System erfüllt alle geforderten Funktionen und bietet eine solide Grundlage für eine Bibliotheksverwaltung. Die Architektur ist sauber getrennt und erweiterbar.

\newpage

% ============================================================================
% ANHANG
% ============================================================================
\appendix
\chapter{Anhang}

\section{Installationsanleitung}

\subsection{Voraussetzungen}
\begin{itemize}
    \item Docker Desktop
    \item Java 17 (JDK)
    \item Node.js 18+
    \item Git
\end{itemize}

\subsection{Setup}
\begin{lstlisting}[language=bash, caption=Installation]
# Repository klonen
git clone https://github.com/hereqi/Probe-IPA-Buecherrei.git
cd Probe-IPA-Buecherrei

# Docker Container starten (MariaDB + Keycloak)
cd docker
docker-compose up -d

# Warten bis Keycloak gestartet ist (~30 Sekunden)
# Keycloak UI: http://localhost:8180 (admin/admin)

# Backend starten
cd ../backend
mvn spring-boot:run

# Frontend starten (neues Terminal)
cd ../frontend
npm install
npm run dev
\end{lstlisting}

\subsection{URLs}
\begin{table}[H]
\centering
\begin{tabular}{l l}
\toprule
\textbf{Service} & \textbf{URL} \\
\midrule
Frontend & \url{http://localhost:5173} \\
Backend API & \url{http://localhost:8080/api} \\
Keycloak Admin & \url{http://localhost:8180} \\
\bottomrule
\end{tabular}
\end{table}

\subsection{Test-Accounts}
\begin{table}[H]
\centering
\begin{tabular}{l l l l}
\toprule
\textbf{Rolle} & \textbf{Benutzername} & \textbf{Passwort} & \textbf{Berechtigung} \\
\midrule
Mitarbeiter & mitarbeiter & test123 & Vollzugriff \\
Kunde & kunde & test123 & Nur eigene Ausleihen \\
\bottomrule
\end{tabular}
\caption{Keycloak Test-Accounts}
\end{table}

\newpage

\section{Quellenverzeichnis}
\begin{itemize}
    \item React Dokumentation: \url{https://react.dev}
    \item Spring Boot Dokumentation: \url{https://spring.io/projects/spring-boot}
    \item Keycloak Dokumentation: \url{https://www.keycloak.org/documentation}
    \item Shadcn UI: \url{https://ui.shadcn.com}
    \item Tailwind CSS: \url{https://tailwindcss.com}
    \item MariaDB: \url{https://mariadb.org}
    \item Docker: \url{https://docs.docker.com}
\end{itemize}

\section{Glossar}
\begin{table}[H]
\begin{tabularx}{\textwidth}{l X}
\toprule
\textbf{Begriff} & \textbf{Erklärung} \\
\midrule
API & Application Programming Interface \\
CORS & Cross-Origin Resource Sharing \\
DTO & Data Transfer Object \\
JPA & Java Persistence API \\
JWT & JSON Web Token \\
OIDC & OpenID Connect \\
ORM & Object-Relational Mapping \\
REST & Representational State Transfer \\
SPA & Single Page Application \\
\bottomrule
\end{tabularx}
\end{table}

\end{document}
